\subsection{传递环节比较}

\paragraph{P 环节}
在单位阶跃输入下，系统会立即响应输入信号的变化，并无延迟地达到其稳态终值。
在正弦输入下，输出信号被放大但没有相位偏移，因为该环节只进行比例放大。

\paragraph{I 环节}
在单位阶跃输入下，系统输出量线性上升，因为它对输入信号随时间进行积分。
这种线性斜率是积分系统的典型特征。

\paragraph{PT1 环节}
在单位阶跃输入下，系统经过一定过渡时间后达到稳态终值。
它以指数形式接近该终值，其中延迟由时间常数 $T$ 决定。
在正弦输入下，与 P 环节不同，PT1 环节会出现相位偏移和幅值减小，这是具有储能行为系统的典型特征。

\paragraph{PT2 环节}
在单位阶跃输入下，PT2 环节表现出比一阶系统更复杂的过渡过程。
它接近稳态终值的速度取决于固有频率和阻尼系数，并且根据阻尼大小可能出现超调。
在正弦输入下，与 PT1 环节一样会出现相位偏移和幅值减小。
由于系统阶次更高，振荡和阻尼效应比 PT1 环节更加明显。

\subsection{阶跃响应表示}

当固有频率 $\omega$ 保持不变时，阻尼系数 $D$ 的影响如下：
\begin{itemize}
  \item $D = 0$（无阻尼）：系统出现幅值恒定的持续振荡，不收敛到稳态值。
  \item $0 < D < 1$（弱阻尼）：系统发生超调，但幅值随时间逐渐减小，直到达到稳态终值。
  \item $D = 1$（非周期临界情况）：系统在无超调条件下较快接近稳态值，且不发生振荡。
  \item $D > 1$（强阻尼）：系统不发生振荡，但响应非常迟缓，需要明显更长时间才能达到稳态终值。
\end{itemize}

当阻尼系数 $D$ 保持不变时，固有频率 $\omega$ 的影响如下：
\begin{itemize}
  \item $\omega$ 越大，系统响应越快，并在更短时间内达到平衡状态。
  \item $\omega$ 减小时，系统变得更慢、更迟缓；若存在振荡，振荡频率也会降低。
\end{itemize}

综上，PT2 系统的过渡行为主要可以通过参数 $\omega$ 和 $D$ 表征。
固有频率 $\omega$ 主要决定响应速度，而阻尼系数 $D$ 决定稳定性以及振荡或超调倾向。

\subsection{Bode 图}

Bode 图显示了 PT2 系统的典型频率特性。
当阻尼较小时，幅值曲线在固有频率附近可能出现共振峰，随后在高频范围内下降。
相位曲线从 $0^\circ$ 开始，并随频率升高逐渐减小。
当达到固有频率 $\omega_n$ 时，相位约为 $-90^\circ$；在很高频率下，最终接近 $-180^\circ$。

\subsubsection{第一组}

第一组在固定固有频率下改变阻尼 $D$。
较小的阻尼会导致更高的峰值（共振峰）以及在设定频率附近更明显的振荡。
随着阻尼增大，响应曲线更加平滑，峰值减小。
同时，相位过渡更加平缓，共振附近的突变减弱。

\subsubsection{第二组}

第二组在固定阻尼系数 $D$ 下改变固有频率 $\omega$。
随着 $\omega$ 增大，系统的共振频率在图中向右移动。
因此幅值峰值也向右移动，其高度和形状仍由固定的阻尼系数决定。
与幅值类似，随着固有频率增大，相位变化的拐点也向较高频率移动。

\subsection{零极点图}

\subsubsection{结构与特征}

零极点图在复 $s$ 平面（$s = \sigma + j\omega_d$）中几何表示系统传递函数的奇异点。
经典 PT2 环节恰好具有两个极点（分母多项式的零点），而分子中没有零点。
这些极点的几何位置从根本上决定系统的稳定性和过渡行为：
\begin{itemize}
  \item 稳定性：由于极点位于左半平面（$\sigma < 0$），系统是绝对稳定的。
  \item 实部（$\sigma = -D\omega_n$）：它决定过渡过程的衰减速度。实部越负，振荡衰减越快。
  \item 虚部（$\omega_d = \omega_n \sqrt{1 - D^2}$）：它表示阻尼固有频率，也就是时域中的实际振荡频率。
\end{itemize}

\subsubsection{第一组：改变阻尼系数}

第一实验组系统改变阻尼系数 $D$，同时将固有频率 $\omega_n$ 固定为常数。
实验数据表现出如下极点位置变化：
\begin{itemize}
  \item 当 $0 < D < 1$（欠阻尼）时，极点为共轭复数。改变 $D$ 时，它们沿以原点为中心的圆弧移动，该圆弧半径等于恒定的固有频率 $\omega_n$。
  \item 当 $D = 0$（无阻尼情况）时，极点为纯虚数，位于竖直轴上（$s = \pm j\omega_n$）。
  \item 当 $D = 1$（非周期临界情况）时，两个极点在负实轴上重合，并在 $s = -\omega_n$ 处形成一个二重实极点。
  \item 当 $D > 1$（过阻尼）时，极点在实轴上分离，并作为两个不同的实极点继续移动。
\end{itemize}

阻尼系数 $D$ 与极点射线相对于负实轴的夹角 $\varphi$ 直接相关（$\cos\varphi = D$）。
减小 $D$ 会使极点更靠近虚轴：虚部增大，实部绝对值减小。
这会在时域中导致系统阻尼减小，表现为明显超调和显著延长的调节时间。

\subsubsection{第二组：改变固有频率}

第二实验组改变固有频率 $\omega_n$，同时保持阻尼系数 $D$ 不变。
由于实部和虚部之间的比例由恒定的 $D$ 决定，极点沿一条以固定角度 $\varphi$ 从原点出发的直线移动。
随着 $\omega_n$ 增大，极点的实部 $\sigma$ 和虚部 $\omega_d$ 均按比例增大，因此极点进一步向外并更深入左半平面。

增大固有频率 $\omega_n$ 会使极点实部进一步向负方向移动。
因此系统响应速度明显提高，调节时间显著缩短。
尽管振荡频率更高，相对阻尼仍保持不变。
只要极点仍位于左半平面，系统就是绝对稳定的。

\subsubsection{总结}

该图从数学上表示 PT2 系统。
它表明，极点到原点的径向距离描述系统的速度 $\omega_n$，而相对于实轴的角度分量决定振荡倾向 $D$。
因此，Bode 图或阶跃响应中的任何变化都可以直接由这两个系统极点的位置变化推导出来。
